Les points, 
variables et de manière plus générale, 
tout élément nommé dans la partie qui suit se retrouve dans le schéma du prisme ci-dessus.~(\ref{fig:schema_prisme}) \\
Il existe quatre lois du prisme communément admises dans le modèle de l'optique géométrique, présentées telles que suit :

\paragraph{Lois de Snell-Descartes en $I$} 
L'expression des lois de la réflexion au point I correspond à :
\[n_1 \cdot \sin(i) = n_2 \cdot \sin(r)\] 
Puisque dans notre contexte, $n_1 = n_{air} = 1$, que $n_2 = n_{{prisme}} = n$, l'expression devient : 
\begin{equation}
    \sin(i) = n \cdot \sin(r)
\end{equation}


\paragraph{Lois de Snell-Descartes en $I'$} 
Concernant le point $I'$, on retrouve de manière parfaitement analogue l'expression : 
\begin{equation}
    \sin(i') = n \cdot \sin(r)
\end{equation}


